% !TEX root = ../macro.tex
\clearpage
\part{Open Economy}\label{part:open-economy}

An open economy can trade goods and claims on future income. This changes the
connection between domestic saving and investment, introduces relative prices
between countries, and allows shocks to cross borders. We extend the earlier
models before comparing policy choices in Part IV.

\clearpage
\plannedsection{International Accounts and Intertemporal Trade}\label{sec:international-accounts}

\subsection{Trade balances, current accounts, and external wealth}
Derive the links between national saving, investment, and the current account.
Distinguish transaction flows from valuation changes in the international
investment position.

\subsection{Consumption smoothing and international borrowing}
Start with a small open endowment economy. Derive its intertemporal budget
constraint and explain how expected income affects the current account.

\subsection{Investment and external adjustment}
Add capital accumulation and examine how borrowing finances investment.
State borrowing and solvency conditions before discussing external adjustment.

\plannedsection{Trade and Relative Prices}\label{sec:trade}

\subsection{Comparative advantage, trade gains, and trade costs}
Use a simple trade model to explain specialization and gains from exchange.
Introduce tariffs and other trade costs to connect trade policy to allocation.

\subsection{Terms of trade and the real exchange rate}
Define both relative prices with explicit conventions. Trace how their changes
affect purchasing power, production, and external balances.

\subsection{Tradable goods, nontradable goods, and adjustment}
Introduce the distinction between tradable and nontradable production.
Explain how expenditure and resources move between sectors during adjustment.

\clearpage
\plannedsection{Open Economy Fluctuations and Exchange Rates}\label{sec:open-fluctuations}

\subsection{The small open economy RBC model}
Extend the real benchmark to international borrowing and investment.
Compare productivity, world-interest-rate, and terms-of-trade shocks.

\subsection{Capital flows, interest parity, and risk}
Derive the relevant asset-pricing conditions and distinguish covered parity
from an uncovered expected-return condition. Introduce risk premia and
foreign-currency balance-sheet exposure.

\subsection{Nominal rigidities and exchange-rate regimes}
Introduce an open economy New Keynesian benchmark. Compare the adjustment
mechanisms under fixed and flexible exchange rates, and state the constraints
created by capital mobility.
